\documentclass[11pt]{article}
\usepackage[margin=0.9in]{geometry}
\usepackage{amsmath,amssymb}
\usepackage{booktabs}
\usepackage{array}
\usepackage{caption}
\usepackage[protrusion=true,expansion=false]{microtype}
\usepackage{parskip}

\title{PML Assignment 2: Self-Supervised Denoising for Noise-Level Selection}
\author{Brian Butler}
\date{March 2026}

\begin{document}
\maketitle

\section{Overview}

We study self-supervised denoising as a pretraining task for a downstream
classification problem on the Iris dataset ($n{=}150$, $d{=}4$, $K{=}3$). Each
input is corrupted by additive isotropic Gaussian noise,
\[
\tilde{x} = x + \varepsilon, \qquad \varepsilon \sim \mathcal{N}(0, \sigma^2 I),
\]
and a Ridge regressor $g_\sigma$ is trained to reconstruct $x$ from $\tilde{x}$
by minimising $\mathbb{E}\,\lVert x - g_\sigma(\tilde{x}) \rVert^2$. Here
$\sigma$ is the \textbf{standard deviation of the corruption noise}, measured
in units of standardised features, and it is the single hyperparameter that
controls the difficulty of the pretext task.

\textbf{Why we do this.} The denoiser $g_\sigma$ never sees any labels, yet in
solving the reconstruction task it must learn the correlation structure of the
features and project them onto a lower-dimensional data manifold. The resulting
$g_\sigma$ is then reused as a \emph{feature transformer} for a classifier
trained on only a handful of labels. The role of $\sigma$ is to tune how
aggressive this implicit regularisation is: if $\sigma$ is too small,
$g_\sigma$ collapses to the identity and learns nothing; if $\sigma$ is too
large, the target becomes unrecoverable and $g_\sigma$ collapses toward the
feature mean. A moderate $\sigma$ sits in between and should, in principle,
yield a useful denoised representation. The central question of the assignment
is how to pick that $\sigma$.

We compare two strategies: (1)~$\sigma_{\text{proxy}}$ --- chosen without
labels by minimising proxy reconstruction MSE; (2)~$\sigma_{\text{down}}$ ---
chosen with few labels ($m{=}5$/class) by maximising test accuracy. Features
are standardised using training-split statistics only to prevent leakage.
Ridge regression ($\alpha{=}1$) is the denoiser; multinomial logistic
regression is the classifier. In Part~B, $g_\sigma$ is used as a consistent
feature transformer: clean inputs pass through $g_\sigma$ for both labeled and
test sets, avoiding any train/test mismatch.

\section{Results}

\subsection{Part A -- Proxy-Task (No Labels)}
Features are split 75/25 (seed 42). For each $\sigma$, noise is added, a Ridge
denoiser is fitted, and MSE is evaluated on the held-out set.

\begin{table}[h]
\centering
\begin{tabular}{cc}
\toprule
$\sigma$ & Proxy-Val MSE \\
\midrule
0.0 & 0.000685 \\
0.1 & 0.008499 \\
0.2 & 0.027897 \\
0.3 & 0.054187 \\
0.4 & 0.085045 \\
0.6 & 0.155183 \\
0.8 & 0.230314 \\
\bottomrule
\end{tabular}
\caption{$\sigma$ vs.\ Proxy-Val MSE.}
\end{table}

Rule: smallest MSE among $\sigma > 0$. \quad
\textbf{Result: $\sigma_{\text{proxy}} = 0.1$.}
$\sigma{=}0$ is excluded (trivial identity). Among $\sigma > 0$, the lowest
noise yields the best reconstruction.

\subsection{Part B -- Downstream (Few Labels)}
Data split 50/50 (stratified, seed 42); $m{=}5$ labeled per class. Clean
features are passed through $g_\sigma$ for both labeled training and test sets.

\begin{table}[h]
\centering
\begin{tabular}{cc}
\toprule
$\sigma$ & Test Accuracy \\
\midrule
0.0 & 0.8533 \\
0.1 & 0.8533 \\
0.2 & 0.8533 \\
0.3 & 0.8533 \\
0.4 & 0.8667 \\
0.6 & 0.8667 \\
0.8 & 0.8667 \\
\bottomrule
\end{tabular}
\caption{$\sigma$ vs.\ Test Accuracy ($m{=}5$/class).}
\end{table}

Rule: highest test accuracy; ties broken by smallest $\sigma$. \quad
\textbf{Result: $\sigma_{\text{down}} = 0.4$.}

\subsection{Multi-Seed Variance (Part B, 10 seeds)}
The single-seed results above show many ties (accuracy differences
$\leq 0.013$). To assess stability, Part~B is repeated over 10 independently
sampled labeled sets. $\mathrm{SEM} = \mathrm{std}/\sqrt{n}$ is reported
alongside the mean.

\begin{table}[h]
\centering
\begin{tabular}{cccc}
\toprule
$\sigma$ & Mean Acc & Std & SEM ($n{=}10$) \\
\midrule
0.0 & 0.8893 & 0.0408 & 0.0129 \\
0.1 & 0.8867 & 0.0441 & 0.0139 \\
0.2 & 0.8827 & 0.0498 & 0.0158 \\
0.3 & 0.8787 & 0.0433 & 0.0137 \\
0.4 & 0.8693 & 0.0478 & 0.0151 \\
0.6 & 0.8560 & 0.0365 & 0.0116 \\
0.8 & 0.8573 & 0.0361 & 0.0114 \\
\bottomrule
\end{tabular}
\caption{$\sigma$ vs.\ Mean / Std / SEM Test Accuracy (10 seeds).}
\end{table}

At $n{=}10$, $\sigma{=}0$ yields the highest mean accuracy (0.8893). But every
SEM is $\geq 0.011$, and the top four $\sigma$ values fall within one SEM of
each other. Before drawing conclusions we therefore need to ask whether 10
replicates are enough.

\subsection{Are 10 Replicates Sufficient?}
To quantify how much of the $n{=}10$ ranking is real signal versus sampling
noise, we extend the experiment to a reference run of $n{=}200$ independent
labeled-set seeds, then evaluate three diagnostics at
$n \in \{5, 10, 20, 50, 100, 200\}$.

\textbf{(a) Standard error of the mean.} SEM shrinks as $1/\sqrt{n}$. At
$n{=}10$ the SEM on each $\sigma$-cell is $0.011$--$0.016$, which is
\emph{larger} than the mean-accuracy gaps between adjacent $\sigma$ values
($0.003$--$0.013$). That alone tells us the $n{=}10$ ranking of the top few
$\sigma$ cannot be trusted.

\begin{table}[h]
\centering
\small
\begin{tabular}{ccccccccc}
\toprule
$n$ & $\sigma{=}0$ & $\sigma{=}0.1$ & $\sigma{=}0.2$ & $\sigma{=}0.3$ & $\sigma{=}0.4$ & $\sigma{=}0.6$ & $\sigma{=}0.8$ \\
\midrule
5   & 0.0244 & 0.0265 & 0.0294 & 0.0258 & 0.0269 & 0.0232 & 0.0225 \\
10  & 0.0129 & 0.0139 & 0.0158 & 0.0137 & 0.0151 & 0.0116 & 0.0114 \\
20  & 0.0084 & 0.0087 & 0.0102 & 0.0097 & 0.0099 & 0.0081 & 0.0085 \\
50  & 0.0078 & 0.0077 & 0.0078 & 0.0072 & 0.0070 & 0.0060 & 0.0055 \\
100 & 0.0052 & 0.0051 & 0.0052 & 0.0049 & 0.0048 & 0.0041 & 0.0037 \\
200 & 0.0036 & 0.0035 & 0.0036 & 0.0034 & 0.0032 & 0.0028 & 0.0025 \\
\bottomrule
\end{tabular}
\caption{SEM of the mean test accuracy at each $n$.}
\end{table}

\textbf{(b) Decision stability.} We bootstrap 1000 subsamples of size $n$ from
the 200 available seeds and ask how often $\arg\max_\sigma$ of the subsample
mean matches the $n{=}200$ reference (which is $\sigma{=}0.1$). At $n{=}10$
the two agree only 22\% of the time, barely better than chance
($1/7 \approx 14\%$). Even at $n{=}200$ the match rate is only 53\%, because
the top three $\sigma$ values ($0,\,0.1,\,0.2$) have overlapping means and are
genuinely indistinguishable on this dataset.

\begin{table}[h]
\centering
\begin{tabular}{ccc}
\toprule
$n$ & $P(\arg\max \text{ matches ref})$ & Avg 95\% CI half-width \\
\midrule
5   & 0.149 & 0.0500 \\
10  & 0.223 & 0.0264 \\
20  & 0.238 & 0.0177 \\
50  & 0.319 & 0.0137 \\
100 & 0.355 & 0.0093 \\
200 & 0.528 & 0.0063 \\
\bottomrule
\end{tabular}
\caption{Stability of the $\arg\max$-$\sigma$ decision (reference: $n{=}200$, $\sigma{=}0.1$).}
\end{table}

\textbf{(c) Paired comparison.} When seed-level noise is the dominant variance
component, pairing across seeds is much more powerful than comparing means. A
paired $t$-test for $\sigma{=}0$ vs.\ $\sigma{=}0.4$ (the two headline values
from the tables above) is significant even at $n{=}10$ (mean
$\Delta{=}{+}0.020$, $t{=}3.00$, $p{=}0.015$). So while 10 seeds cannot pick
the single best $\sigma$, they are enough to reject the single-seed claim
that $\sigma{=}0.4$ beats $\sigma{=}0$.

\begin{table}[h]
\centering
\begin{tabular}{ccccc}
\toprule
$n$ & mean diff & SE & $t$ & $p$ \\
\midrule
5   & $+0.0267$ & 0.0094 & $+2.83$ & 0.047 \\
10  & $+0.0200$ & 0.0067 & $+3.00$ & 0.015 \\
20  & $+0.0133$ & 0.0053 & $+2.52$ & 0.021 \\
50  & $+0.0128$ & 0.0032 & $+3.99$ & 0.0002 \\
100 & $+0.0084$ & 0.0026 & $+3.22$ & 0.002 \\
200 & $+0.0092$ & 0.0019 & $+4.96$ & $<0.001$ \\
\bottomrule
\end{tabular}
\caption{Paired $t$-test, $\sigma{=}0$ vs.\ $\sigma{=}0.4$ (per-seed pairing).}
\end{table}

\textbf{Verdict.} Ten replicates are \emph{not} sufficient to identify the
best $\sigma$: the 95\% CI on each mean is $\approx \pm 0.026$, wider than the
spread of the top four $\sigma$ values, and the $\arg\max$ choice at $n{=}10$
is unstable (matches the reference only 22\% of the time). Ten replicates
\emph{are} sufficient to detect the broad trend that low $\sigma$ outperforms
high $\sigma$ (paired $t$-test, $p{=}0.015$). For a reliable point estimate,
$n \geq 50$ is advisable ($|\text{mean} - \text{reference}| < 0.005$). In
general, whenever candidate $\sigma$ values produce mean accuracies separated
by less than a SEM, the honest reporting format is mean $\pm$ SEM with a
pairwise significance test, not a single $\arg\max$.

\section{Discussion}

\subsection{Why Do the Two Strategies Disagree?}
The proxy task optimises reconstruction fidelity; the downstream task
optimises classification accuracy. These are fundamentally different
objectives. $\sigma_{\text{proxy}}{=}0.1$ minimises reconstruction error among
non-trivial noise levels: small noise forces the denoiser to learn data
structure while keeping distortion low. The single-seed
$\sigma_{\text{down}}{=}0.4$ arises because the denoiser applies a mild
low-pass filter that happened to help on one particular labeled subset. The
multi-seed (and the $n{=}200$ reference) analysis in \S2.4 shows that on
average $\sigma \approx 0.1$ is in fact best, and that $\sigma{=}0.4$ is
significantly worse. So the real divergence is modest
($\sigma \approx 0.1$ vs.\ $\sigma{=}0$ being essentially tied), reflecting
that on Iris the raw features are already informative enough that denoising
adds little benefit.

\subsection{Is There a ``Sweet Spot''?}
For reconstruction, MSE increases monotonically with $\sigma$; the sweet spot
is the smallest positive $\sigma$ ($0.1$). For classification, the $n{=}200$
means decrease gradually from $0.8683$ at $\sigma{=}0$ to $0.8521$ at
$\sigma{=}0.8$, with the top three $\sigma$ values statistically
indistinguishable. The ``sweet spot'' is therefore broad and centred near
$\sigma{=}0$. On higher-dimensional or noisier datasets one would expect a
clear interior peak where moderate denoising regularises representations and
improves generalisation.

\subsection{Proxy-Based vs.\ Downstream Tuning}
\textbf{Proxy-based ($\sigma_{\text{proxy}}$):} requires no labels; scalable
and stable. However, the proxy objective may not align with downstream needs.

\textbf{Downstream ($\sigma_{\text{down}}$):} directly optimises the target
metric, but requires labeled data and --- as \S2.4 shows --- is highly
sensitive to sampling variance with small budgets. A single labeled-set draw
can easily mislead; even ten cannot reliably select a winner. Multi-seed
averaging with reported SEM, combined with paired significance tests, is
essential.

\textbf{When to prefer each:} use proxy tuning when labels are unavailable or
when the labeled budget is too small for the downstream signal to overcome
seed-level variance; use downstream tuning when even a small labeled set
exists \emph{and} enough seeds can be run ($n \geq 50$) to stabilise the
estimate.

\subsection{Reproducibility}
Three sources of randomness are fixed: (1)~\texttt{random\_state=42} in
\texttt{train\_test\_split} ensures identical partitions;
(2)~\texttt{RandomState(42)} for labeled-set sampling guarantees the same
$m{=}5$ examples per class in the single-seed run; (3)~distinct seed offsets
(\texttt{SEED+1}, \texttt{SEED+100+i}) provide independent noise and
labeled-set realisations while remaining reproducible. Parts~A and B use
separate train/test splits (75/25 vs.\ 50/50) since they are independent
experiments; within each part, splits and labeled budgets are held constant
across all $\sigma$ values.

\end{document}
